% Generated draft: requires mathematical and source review.
\documentclass[11pt]{article}
\usepackage[margin=25mm]{geometry}
\usepackage{fontspec}
\setmainfont{Latin Modern Roman}
\usepackage{amsmath,amssymb,mathtools,unicode-math}
\setmathfont{Latin Modern Math}
\usepackage{enumitem,needspace,xurl,hyperref}
\setlength{\parindent}{0pt}
\setlength{\parskip}{5pt}
\setlength{\emergencystretch}{3em}
\newcommand{\N}{\mathbb{N}}
\newcommand{\Z}{\mathbb{Z}}
\newcommand{\Q}{\mathbb{Q}}
\newcommand{\R}{\mathbb{R}}
\newcommand{\C}{\mathbb{C}}
\newcommand{\F}{\mathbb{F}}
\newcommand{\A}{\mathbb{A}}
\newcommand{\PP}{\mathbb{P}}
\newcommand{\E}{\mathbb{E}}
\newcommand{\eps}{\varepsilon}
\newcommand{\dd}{\,\mathrm{d}}
\begin{document}
\section*{UnsolvedMath ID: 2306041}
\textbf{Source problem code:} catalogue entry
\subsection{Definitions and mathematical statement}
Let \(S\) be the normalized univalent holomorphic functions on \(\mathbb D\): \(f(0)=0\), \(f'(0)=1\). Define
\[
K(\alpha)=\left\{f\in S:\operatorname{Re}\left(1+\frac{zf''(z)}{f'(z)}\right)\ge\alpha\quad(z\in\mathbb D)\right\},
\]
\[
S^*(\alpha)=\left\{f\in S:\operatorname{Re}\frac{zf'(z)}{f(z)}\ge\alpha\quad(z\in\mathbb D)\right\}.
\]
The quotient at \(0\) is interpreted by its analytic extension, equal to \(1\). The source omits the range of \(\alpha\); the customary \(0\le\alpha<1\) range is a possible intended domain, not an unqualified source hypothesis. For \(0<r<1\), part (a) asks to prove
\[
\min_{f\in K(\alpha)}\min_{|z|=r}\left|\frac{zf'(z)}{f(z)}\right|
=\min_{f\in K(\alpha)}\min_{|z|=r}\operatorname{Re}\frac{zf'(z)}{f(z)}.
\]
The source uses minima; existence of extremizers or an intended infimum convention requires checking.

Part (b) asks to prove specified starlikeness orders for
\[
h_\alpha(z)=\frac{1-(1-z)^{2\alpha-1}}{2\alpha-1}\quad(\alpha\ne\tfrac12),
\qquad h_{1/2}(z)=-\log(1-z),
\]
where the analytic branch has \(\log(1)=0\). The respective claimed orders are printed as
\[
\frac{4^\alpha(2\alpha-1)}{4-4.2^\alpha},\qquad \frac1{\log4}.
\]
The denominator \(4-4.2^\alpha\) is reproduced literally and is ambiguous. Interpreting the dot as multiplication gives \(4-4\cdot2^\alpha\), which vanishes at \(\alpha=0\) and does not yield the stated half-order limit. No corrected formula is guessed. Once that expression and parameter range are fixed, the requested membership is \(h_\alpha\in S^*(\beta_\alpha)\).

The source says either part, together with I. S. Jack's work, would solve F. R. Keogh's problem, printed as
\[
\max_{f\in K(\alpha)}\{\beta:f\in S^*(\beta)\}.
\]
This is not a precise scalar optimization as written. Moreover, maximizing individual starlikeness order is different from finding the largest \(\beta\) with \(K(\alpha)\subseteq S^*(\beta)\); in the customary range the identity \(f(z)=z\) already has individual order \(1\). The intended uniform-inclusion constant may instead be
\[
\sup\{\beta:K(\alpha)\subseteq S^*(\beta)\},
\]
but that is explicitly a possible interpretation, not a replacement of the source. The supplied OPEN-TRIAGE summary reports open status in the 2018 edition and no definitive resolution located.
\subsection{Short English statement}
Determine extremal starlikeness information for convex functions of order \(\alpha\). The record includes two proposed routes, but its parameter range, explicit extremal-order formula, and final maximum-versus-uniform-inclusion target need checking.
\subsection{Sources}
\begin{itemize}
\item Local UnsolvedMath registry, permanent ID \(2306041\), exact source code \texttt{catalogue entry}, ``Research Problems in Function Theory --- Problem 6.41.''
\item Hayman, \emph{Research Problems in Function Theory} (2018), Problem 6.41, attributed to D. Benjamin and referencing F. R. Keogh's problem; \url{https://arxiv.org/abs/1809.07200}. Source-TeX extraction; recorded access 2026-07-29.
\item Source references Goel for part (a), MacGregor for part (b), and I. S. Jack for the proposed implication, without complete citations. The supplied OPEN-TRIAGE summary gives no separate resolution reference.
\end{itemize}
\end{document}
